\documentclass[11pt,a4paper]{article}

\usepackage[T1]{fontenc}
\usepackage[utf8]{inputenc}
\usepackage{lmodern}
\usepackage{microtype}
\usepackage[a4paper,margin=1in]{geometry}
\usepackage{amsmath}
\usepackage{amssymb}
\usepackage{amsthm}
\usepackage{mathtools}
\usepackage{booktabs}
\usepackage{xcolor}
\usepackage[colorlinks=true,linkcolor=blue!55!black,citecolor=blue!55!black,
            urlcolor=blue!55!black,breaklinks=true]{hyperref}

\newtheorem{theorem}{Theorem}
\newtheorem{lemma}{Lemma}
\newtheorem{corollary}{Corollary}
\theoremstyle{remark}
\newtheorem{remark}{Remark}

\newcommand{\id}{\mathrm{id}}
\newcommand{\Tr}{\operatorname{Tr}}
\newcommand{\Ad}{\operatorname{Ad}}
\newcommand{\cD}{\mathcal{D}}
\newcommand{\cE}{\mathcal{E}}
\newcommand{\cN}{\mathcal{N}}
\newcommand{\cT}{\mathcal{T}}
\newcommand{\cU}{\mathcal{U}}
\newcommand{\cL}{\mathcal{L}}
\newcommand{\diamondnorm}[1]{\left\lVert #1\right\rVert_{\diamond}}
\newcommand{\tracenorm}[1]{\left\lVert #1\right\rVert_{1}}
\newcommand{\SWAP}{\mathrm{SWAP}}

\hypersetup{
  pdftitle={Exact directional signalling for Ising and partial-SWAP unitary families},
  pdfauthor={Richard Astbury},
  pdfkeywords={quantum signalling; diamond norm; partial SWAP; Ising interaction;
               quantum channels; semidefinite programming}
}

\title{\bfseries Exact directional signalling for Ising and
partial-SWAP unitary families}

\author{Richard Astbury\thanks{Contact:
\href{mailto:richard.astbury@gmail.com}{richard.astbury@gmail.com}.
Developed with the assistance of AI tools. The author is solely responsible
for all claims, proofs, code, and errors.}}

\date{Draft of 23 July 2026}

\begin{document}
\maketitle

\begin{abstract}
For a bipartite unitary $U$ on $A\otimes B$, directional signalling from
$B$ to $A$ is the diamond-norm distance
\[
 \delta_A(U)=\inf_{\cE\ {\rm CPTP}}
 \diamondnorm{\Tr_B\!\circ\Ad_U-\cE\circ\Tr_B}.
\]
This quantity was introduced by Barsse, Perinotti, Tosini, and Vaglini.
Exact values were subsequently obtained for the CNOT and SWAP gates, but
continuous interacting families have not been evaluated.  We give two such
evaluations for qubits.  First, for the Ising family
$U_\theta=\exp(-i\theta Z\otimes Z)$, we prove
$\delta_A(U_\theta)=|\sin 2\theta|$ and identify an optimal dephasing channel.
Second, for the partial-SWAP family
$U_\phi=\cos\phi\,I-i\sin\phi\,\SWAP$, unitary covariance reduces the
optimization to one depolarizing parameter $\lambda$.  We evaluate the
fixed-$\lambda$ diamond norm in closed form and thereby reduce the full
problem to a scalar convex minimization.  In the weak-interaction regime
$0\leq\sin\phi\leq 1/3$, the result simplifies to
$\delta_A(U_\phi)=2\sin\phi$; beyond this threshold the optimal effective
channel moves continuously away from the identity and obeys an explicit
algebraic stationarity equation, reaching the known SWAP value $3/2$.
The proofs include matching operational witnesses.  We also distinguish
these diamond-norm results from Hilbert--Schmidt operator-entanglement
measures and provide exact-arithmetic computational certificates.
\end{abstract}

\noindent\textbf{Keywords:} directional signalling; diamond norm;
partial SWAP; Ising interaction; covariant quantum channels.

\section{Introduction}

An interacting microscopic evolution need not induce autonomous dynamics on
a subsystem.  Given a unitary channel $\Ad_U$ on $A\otimes B$, one may ask how
well the reduced output on $A$ can be simulated from the reduced input on $A$
alone.  The operational answer is
\begin{equation}
 \delta_A(U)
 =\inf_{\cE\in\mathrm{CPTP}(A)}
   \diamondnorm{\Tr_B\circ\Ad_U-\cE\circ\Tr_B}.
 \label{eq:defect}
\end{equation}
The optimization includes arbitrary $AB$ correlations and arbitrary reference
systems through the diamond norm.

Equation~\eqref{eq:defect} is not a new measure.  After exchanging subsystem
labels, it is exactly the directional-signalling quantity introduced in
Ref.~\cite{barsse2024}.  That work bounded the CNOT value; the follow-up
Ref.~\cite{barsse2025} computed exact CNOT and SWAP values, including parallel
and asymptotic uses.  Earlier work characterized causal and semicausal
operations structurally~\cite{beckman2001,eggeling2002}, while continuity of
Stinespring dilations supplies general non-sharp bounds~\cite{ksw2008}.

The present contribution is deliberately narrower than a new causal measure.
We compute Eq.~\eqref{eq:defect} sharply for two continuous unitary families.
The Ising result extends the CNOT endpoint.  The partial-SWAP result covers an
interaction used in homogenization and collision models
\cite{ziman2002,scarani2002} and reveals a transition in the optimal effective
channel at $\sin\phi=1/3$.  A dedicated search found no previous evaluation of
Eq.~\eqref{eq:defect} for either continuous family.

The distinction from operator entanglement is important.  Hilbert--Schmidt
distances and the linear entropy of a vectorized unitary are useful algebraic
diagnostics~\cite{zanardi2001}, but they do not directly give operational
channel distinguishability.  Every result below is stated in diamond norm and
is accompanied by both an upper bound and a matching witness.

\section{Preliminaries}

Let $\cL(H)$ denote the operators on a finite-dimensional Hilbert space $H$.
For a unitary $U$, write $\Ad_U(X)=UXU^\dagger$.  We use the unnormalized Choi
operator
\[
 J(\Phi)=\sum_{j,k}\Phi(|j\rangle\langle k|)
                 \otimes |j\rangle\langle k|.
\]
For Hermiticity-preserving $\Phi:\cL(X)\to\cL(Y)$, the diamond norm admits the
Watrous semidefinite-program characterization~\cite{watrous2009,watrous2018}.
In particular, after the symmetries of an SDP have been imposed, one may use
\begin{equation}
 \diamondnorm{\Phi}
 =\max_{\rho\in\mathrm{D}(X)}
 \tracenorm{(I_Y\otimes\sqrt{\rho})J(\Phi)
                     (I_Y\otimes\sqrt{\rho})}.
 \label{eq:watrous}
\end{equation}
The maximizing $\rho$ specifies a physical witness through any purification.

We use the maps
\[
 \cT=\Tr_B,\qquad \cN_U=\Tr_B\circ\Ad_U,
\]
so that $\delta_A(U)=\inf_\cE\diamondnorm{\cN_U-\cE\cT}$.
Pre- and post-composition by unitary channels are diamond-norm isometries.

\begin{lemma}[Covariant twirling]
\label{lem:twirl}
Suppose a compact group $G$ has unitary representations $H_g$ on $A\otimes B$
and $G_g$ on $A$ such that
\[
 \cN_U\Ad_{H_g}=\Ad_{G_g}\cN_U,\qquad
 \cT\Ad_{H_g}=\Ad_{G_g}\cT.
\]
Then some optimal effective channel in Eq.~\eqref{eq:defect} is covariant under
$G_g$.
\end{lemma}

\begin{proof}
For a candidate $\cE$, define
$\cE_g=\Ad_{G_g^\dagger}\cE\Ad_{G_g}$.  Covariance gives
\[
 \cN_U-\cE_g\cT
 =\Ad_{G_g^\dagger}(\cN_U-\cE\cT)\Ad_{H_g},
\]
so every $\cE_g$ has the same defect as $\cE$.  The Haar average
$\overline{\cE}=\int_G\cE_g\,dg$ is CPTP and covariant.  Convexity of the
diamond norm implies that its defect is no larger.
\end{proof}

\section{The Ising family}

Let
\[
 U_\theta=e^{-i\theta Z\otimes Z}
          =cI-is\,Z\otimes Z,\qquad c=\cos\theta,\quad s=\sin\theta.
\]

\begin{theorem}[Exact Ising signalling]
\label{thm:ising}
For two qubits,
\begin{equation}
 \delta_A(U_\theta)=|\sin 2\theta|=2|cs|.
 \label{eq:ising-result}
\end{equation}
An optimal effective channel is
\[
 \cE_\theta(X)=c^2X+s^2ZXZ.
\]
\end{theorem}

\begin{proof}
Define
\[
 M(X)=\Tr_B[(I\otimes Z)X],\qquad
 L(X)=-\frac{i}{2}[Z,M(X)].
\]
Expanding $U_\theta XU_\theta^\dagger$ and using cyclicity inside the partial
trace gives
\begin{equation}
 \cN_{U_\theta}-\cE_\theta\cT
 =2cs\,L.
 \label{eq:ising-factor}
\end{equation}
The completely bounded trace norm of $M$ is one: its adjoint is the complete
operator-norm isometry $Y\mapsto Y\otimes Z$.  Also
$K\mapsto-i[Z,K]/2$ has completely bounded trace norm at most one by the
triangle inequality for left and right unitary multiplication.  Hence
$\diamondnorm{L}\leq1$, which proves the upper bound $2|cs|$.

For the lower bound, consider
\[
 \rho_0=|+\rangle\langle+|_A\otimes|0\rangle\langle0|_B,\qquad
 \rho_1=|+\rangle\langle+|_A\otimes|1\rangle\langle1|_B.
\]
They have the same $A$ marginal, so every candidate $\cE$ assigns them the
same simulated output.  Their true reduced outputs differ by $2cs\,Y$, whose
trace norm is $4|cs|$.  The triangle inequality therefore forces at least one
of the two simulation errors to be at least $2|cs|$.  This matches the upper
bound.
\end{proof}

At $\theta=\pi/4$, $U_\theta$ is locally equivalent to CZ and CNOT.
Equation~\eqref{eq:ising-result} then reproduces the exact CNOT value one,
already obtained in Ref.~\cite{barsse2025}.  The continuous formula, rather
than that endpoint, is the result of Theorem~\ref{thm:ising}.

\section{The partial-SWAP family}

Consider
\begin{equation}
 U_\phi=cI-isS,\qquad S=\SWAP,\qquad c=\cos\phi,\quad s=\sin\phi,
 \label{eq:partial-swap}
\end{equation}
for $0\leq\phi\leq\pi/2$.  The equal-angle Cartan interaction is this family
up to global phase:
\[
 e^{-it(XX+YY+ZZ)}
 =e^{it}e^{-2itS},\qquad \phi=2t,
\]
because $XX+YY+ZZ=2S-I$.

\subsection{Reduction to depolarizing channels}

The unitary $U_\phi$ commutes with $V\otimes V$ for every $V\in U(2)$.
Lemma~\ref{lem:twirl} therefore permits an optimal effective channel to be
chosen covariant under all unitary conjugations.  Every such qubit channel is
depolarizing:
\begin{equation}
 \cE_\lambda(X)
 =\lambda X+(1-\lambda)\Tr(X)\frac{I}{2},
 \qquad -\frac13\leq\lambda\leq1.
 \label{eq:depol}
\end{equation}
The interval is precisely the CPTP range.

For reference, direct expansion of Eq.~\eqref{eq:partial-swap} gives
\[
 \cN_{U_\phi}(X)
 =c^2\Tr_BX+s^2\Tr_AX
  -ics\,\Tr_B(SX-XS),
\]
where $\Tr_A X$ is relabelled as an operator on the output copy of $A$.

\subsection{The fixed-channel diamond norm}

Set
\begin{equation}
 a=\frac{1-\lambda}{3},\qquad
 B=(1+\lambda)^2-4\lambda c^2.
 \label{eq:aB}
\end{equation}
For the CPTP interval in Eq.~\eqref{eq:depol}, $a\geq0$ and $B\geq0$.

\begin{lemma}[Fixed depolarizing channel]
\label{lem:fixed}
Let
\[
 d_\phi(\lambda)
 =\diamondnorm{\cN_{U_\phi}-\cE_\lambda\cT}.
\]
Then
\begin{equation}
 d_\phi(\lambda)=
 \begin{cases}
  4a, & B\leq4a^2,\\[3pt]
  \displaystyle\frac{B}{\sqrt B-a}, & B>4a^2.
 \end{cases}
 \label{eq:fixed-norm}
\end{equation}
\end{lemma}

\begin{proof}
The full linear diamond-norm SDP may be averaged over the simultaneous
$V\otimes V$ action.  This averaging applies to the complete feasible tuple,
not to the nonlinear expression in Eq.~\eqref{eq:watrous}.  By Schur--Weyl
duality an invariant input density has the form
\[
 \rho=xP_-+yP_+,\qquad x+3y=1,
\]
where $P_-$ is the singlet projector and $P_+=I-P_-$ is the triplet
projector.  Averaging the full SDP shows that an optimum of this form exists.

Put
\[
 K_\rho=(I\otimes\sqrt\rho)
 J(\cN_{U_\phi}-\cE_\lambda\cT)
 (I\otimes\sqrt\rho).
\]
In a singlet/triplet basis, elementary block diagonalization produces four
one-dimensional blocks and two identical two-dimensional blocks.  The former
have eigenvalue $(\lambda-1)y/2$; each latter block has trace
$(1-\lambda)y$ and determinant $-3xyB/4$.  Equivalently,
\begin{equation}
 \det(zI-K_\rho)
 =\left(z+\frac{(1-\lambda)y}{2}\right)^4
  \left(z^2-(1-\lambda)yz-\frac{3xyB}{4}\right)^2.
 \label{eq:charpoly}
\end{equation}
Thus four eigenvalues equal $(\lambda-1)y/2$, and the remaining four are two
copies of
\[
 \frac{(1-\lambda)y
 \pm\sqrt{(1-\lambda)^2y^2+3xyB}}{2}.
\]
The two signs give opposite-sign eigenvalues because $B\geq0$.  With
$u=3y=1-x$, Eq.~\eqref{eq:watrous} becomes
\begin{equation}
 d_\phi(\lambda)
 =\max_{0\leq u\leq1}
 2\left[
 au+\sqrt{a^2u^2+Bu(1-u)}
 \right].
 \label{eq:u-max}
\end{equation}

If $B\leq4a^2$, the maximum is at $u=1$ and equals $4a$.  If $B>4a^2$, the
maximizer and value are
\[
 u_*=\frac{\sqrt B}{2(\sqrt B-a)},\qquad
 d_\phi(\lambda)=\frac{B}{\sqrt B-a}.
\]
These maximizing densities, together with any purification, are matching
operational diamond-norm witnesses.
\end{proof}

\subsection{Optimization and the coupling threshold}

\begin{theorem}[Exact partial-SWAP reduction]
\label{thm:partial-swap}
For the qubit partial-SWAP family,
\begin{equation}
 \delta_A(U_\phi)
 =\min_{-1/3\leq\lambda\leq1}d_\phi(\lambda),
 \label{eq:scalar-min}
\end{equation}
where $d_\phi$ is the explicit function in Eq.~\eqref{eq:fixed-norm}.
Moreover,
\begin{equation}
 0\leq\sin\phi\leq\frac13
 \quad\Longrightarrow\quad
 \lambda_*=1,\qquad
 \delta_A(U_\phi)=2\sin\phi.
 \label{eq:weak-branch}
\end{equation}
For $1/3<\sin\phi<1$, an interior minimizer
$0<\lambda_*<1$ obeys
\begin{equation}
 3(\lambda_*+1-2c^2)
 \left[
 \sqrt{B_*}-\frac{2(1-\lambda_*)}{3}
 \right]
 =B_*,
 \label{eq:stationary}
\end{equation}
with $B_*$ defined by Eq.~\eqref{eq:aB}.  At $\phi=\pi/2$,
$\lambda_*=0$ and $\delta_A(U_{\pi/2})=3/2$.
\end{theorem}

\begin{proof}
Equation~\eqref{eq:scalar-min} follows from the covariant-channel reduction and
Lemma~\ref{lem:fixed}.  The function $d_\phi(\lambda)$ is convex because it is
the diamond norm of a map affine in $\lambda$.  At $\lambda=1$,
$a=0$, $B=4s^2$, and
\[
 d_\phi(1)=2s,\qquad
 \left.\frac{d}{d\lambda}d_\phi(\lambda)\right|_{1^-}
 =s-\frac13.
\]
Convexity therefore makes $\lambda=1$ globally optimal exactly on the weak
branch $s\leq1/3$, proving Eq.~\eqref{eq:weak-branch}.

For $1/3<s<1$, moving left from $\lambda=1$ lowers the objective.  At
$\lambda=0$ its derivative is negative when $c>0$, so a minimizer lies in
$(0,1)$.  Differentiating the second line of Eq.~\eqref{eq:fixed-norm} and
setting the derivative to zero yields Eq.~\eqref{eq:stationary}.  At SWAP,
$c=0$, the solution is $\lambda_*=0$, $a=1/3$, $B=1$, and
$d_{\pi/2}(0)=3/2$, in agreement with Ref.~\cite{barsse2025}.
\end{proof}

\begin{table}[t]
\centering
\caption{Representative points on the partial-SWAP curve.  Interior values
are shown only for orientation; the theorem is Eq.~\eqref{eq:scalar-min}.}
\label{tab:values}
\begin{tabular}{@{}lll@{}}
\toprule
$\phi$ & optimal $\lambda_*$ & $\delta_A(U_\phi)$\\
\midrule
$0$ & $1$ & $0$\\
$\arcsin(1/3)$ & $1$ & $2/3$\\
$\pi/4$ & $\approx0.52358$ & $\approx1.31358$\\
$\pi/2$ & $0$ & $3/2$\\
\bottomrule
\end{tabular}
\end{table}

\section{Discussion}

The two families exhibit different mechanisms.  In the Ising family, one
dephasing channel is optimal for every coupling, and two product inputs already
witness the exact defect.  Along the partial-SWAP family, the optimal effective
channel itself changes: it remains the identity until $\sin\phi=1/3$, then
moves into the depolarizing interval and reaches complete depolarization at
SWAP.  The optimal diamond witness likewise generally uses a reference system,
as encoded by a purification of the singlet/triplet density in
Lemma~\ref{lem:fixed}.

The threshold is operational rather than spectral: it is where the outer
channel optimization first benefits from adding depolarization.  It is not
visible from operator entanglement alone.  This supplies a concrete example in
which a coarse-grained autonomy question differs from a Hilbert--Schmidt
interaction diagnostic.

The result is limited in three important ways.  First, Eq.~\eqref{eq:defect}
is a pre-existing signalling measure, not a newly defined quantity.  Second,
the exact CNOT and SWAP endpoints are prior results
\cite{barsse2025}; the contribution here is the continuous-family evaluation.
Third, a closed formula for a generic two-qubit Cartan interaction
$e^{-i(\alpha XX+\beta YY+\gamma ZZ)}$ remains open.  Joint-Pauli symmetry
reduces that problem to optimization over the Pauli-channel tetrahedron, but
the three leakage sectors interfere in trace norm and the partial-SWAP
$U(2)$ symmetry is absent.

\section{Reproducibility}

The accompanying repository implements all finite-dimensional identities over
the Gaussian rationals $\mathbb{Q}(i)$.  It verifies the Ising residual
factorization, the physical product-state witness, the equal-Cartan/partial-SWAP
identity, exact weak-branch values at the Pythagorean point
$(\cos\phi,\sin\phi)=(40/41,9/41)$, and the SWAP endpoint.  The focused suite is
run from the repository root with
\begin{verbatim}
python3 -m unittest discover -s tests -p 'test_aqc_*.py' -v
\end{verbatim}
The code is a certificate and regression harness; the analytic proofs above do
not depend on floating-point optimization.  The exact harness covers an
interior point of the weak branch and the SWAP boundary of the strong branch;
the generic interior strong-branch result is certified analytically by
Eqs.~\eqref{eq:charpoly}--\eqref{eq:stationary}, not by encoding its generally
algebraic minimizer in the rational backend.

\section{Conclusion}

Directional signalling can be evaluated sharply beyond isolated gates.
For Ising interactions it is $|\sin2\theta|$.  For partial SWAP it is an
explicit scalar convex minimization with a closed weak branch
$2\sin\phi$, a transition at $\sin\phi=1/3$, and the known SWAP endpoint
$3/2$.  The calculation provides a controlled setting in which an optimal
autonomous effective channel changes nonperturbatively with interaction
strength.  Extending the same operational analysis to the full Cartan
three-parameter family is the natural next mathematical problem.

\begin{thebibliography}{99}

\bibitem{barsse2024}
K.~Barsse, P.~Perinotti, A.~Tosini, and L.~Vaglini,
``Causal influence versus signalling for interacting quantum channels,''
\emph{Phys. Rev. Research} \textbf{6}, 043305 (2024),
\href{https://arxiv.org/abs/2309.07771}{arXiv:2309.07771}.

\bibitem{barsse2025}
K.~Barsse, P.~Perinotti, A.~Tosini, and L.~Vaglini,
``Disentangling signalling and causal influence'' (2025),
\href{https://arxiv.org/abs/2505.14120}{arXiv:2505.14120}.

\bibitem{beckman2001}
D.~Beckman, D.~Gottesman, M.~A. Nielsen, and J.~Preskill,
``Causal and localizable quantum operations,''
\emph{Phys. Rev. A} \textbf{64}, 052309 (2001),
\href{https://arxiv.org/abs/quant-ph/0102043}{arXiv:quant-ph/0102043}.

\bibitem{eggeling2002}
T.~Eggeling, D.~Schlingemann, and R.~F. Werner,
``Semicausal operations are semilocalizable,''
\emph{Europhys. Lett.} \textbf{57}, 782 (2002),
\href{https://arxiv.org/abs/quant-ph/0104027}{arXiv:quant-ph/0104027}.

\bibitem{ksw2008}
D.~Kretschmann, D.~Schlingemann, and R.~F. Werner,
``A continuity theorem for Stinespring's dilation,''
\emph{J. Funct. Anal.} \textbf{255}, 1889 (2008),
\href{https://arxiv.org/abs/0710.2495}{arXiv:0710.2495}.

\bibitem{watrous2009}
J.~Watrous, ``Semidefinite programs for completely bounded norms,''
\emph{Theory Comput.} \textbf{5}, 217 (2009).

\bibitem{watrous2018}
J.~Watrous, \emph{The Theory of Quantum Information}
(Cambridge University Press, Cambridge, 2018).

\bibitem{zanardi2001}
P.~Zanardi, ``Entanglement of quantum evolutions,''
\emph{Phys. Rev. A} \textbf{63}, 040304(R) (2001),
\href{https://arxiv.org/abs/quant-ph/0010074}{arXiv:quant-ph/0010074}.

\bibitem{ziman2002}
M.~Ziman, P.~Stelmachovic, V.~Buzek, M.~Hillery, V.~Scarani, and N.~Gisin,
``Quantum homogenization,''
\emph{Phys. Rev. A} \textbf{65}, 042105 (2002),
\href{https://arxiv.org/abs/quant-ph/0110164}{arXiv:quant-ph/0110164}.

\bibitem{scarani2002}
V.~Scarani, M.~Ziman, P.~Stelmachovic, N.~Gisin, and V.~Buzek,
``Thermalizing quantum machines: Dissipation and entanglement,''
\emph{Phys. Rev. Lett.} \textbf{88}, 097905 (2002),
\href{https://arxiv.org/abs/quant-ph/0110088}{arXiv:quant-ph/0110088}.

\end{thebibliography}

\end{document}
